\SourceLeaf{232}%
eine bestimmte Seite wählen, auf der wir uns die Figur liegend denken. Auf der Hilfskugel verwenden wir stets die vom Mittelpunkt abgewandte Außenseite. Auf der gekrümmten Fläche kann man ebenfalls die Außenseite oder die als Außenseite betrachtete Seite wählen, genauer gesagt die Seite, auf der man sich die Normale errichtet denkt. Denn an der Ähnlichkeit der Figuren ändert sich offenbar nichts, wenn man auf der gekrümmten Fläche sowohl die Figur als auch die Normale auf die Gegenseite überträgt, solange das Bild der Figur stets auf derselben Seite der Kugeloberfläche gezeichnet wird.

Das positive oder negative Vorzeichen, das wir dem \textit{Maß} der Krümmung je nach der Lage einer unendlich kleinen Figur zuschreiben, übertragen wir auch auf die Integralkrümmung einer endlichen Figur auf der gekrümmten Fläche. Wenn wir den Gegenstand jedoch in voller Allgemeinheit behandeln wollen, bedarf es einiger Erläuterungen, die wir hier nur kurz berühren. Solange die Figur auf der gekrümmten Fläche so beschaffen ist, dass den einzelnen Punkten in ihrem Inneren \textit{verschiedene} Punkte auf der Kugeloberfläche entsprechen, bedarf die Definition keiner weiteren Erklärung. Ist diese Bedingung nicht erfüllt, müssen bestimmte Teile der Figur auf der Kugeloberfläche zweimal oder mehrfach mitgezählt werden. Je nachdem, ob ihre Lage ähnlich oder umgekehrt ist, können sie sich dabei addieren oder aufheben. In diesem Fall denkt man sich die Figur auf der gekrümmten Fläche am einfachsten in Teile zerlegt, die jeder für sich die Bedingung erfüllen. Jedem Teil ordnet man seine Integralkrümmung zu, deren Betrag durch den Flächeninhalt der entsprechenden Figur auf der Kugeloberfläche und deren Vorzeichen durch deren Lage bestimmt wird. Schließlich ordnet man der ganzen Figur die Integralkrümmung zu, die sich durch Addition der Integralkrümmungen der einzelnen Teile ergibt. Allgemein ist also die Integralkrümmung einer Figur \(=\int k\,\mathrm{d}\sigma\), wobei \(\mathrm{d}\sigma\) das Flächenelement der Figur und \(k\) das Krümmungsmaß in einem beliebigen Punkt bezeichnet. Für die geometrische Darstellung dieses Integrals kommt es hauptsächlich auf Folgendes an. Dem Rand der Figur auf der gekrümmten Fläche entspricht unter der Einschränkung von Artikel 3 stets eine in sich zurücklaufende Linie auf der Kugeloberfläche. Schneidet sich diese Linie nirgends selbst, so teilt sie die gesamte Kugeloberfläche in zwei Teile, von denen einer der Figur auf der gekrümmten Fläche entspricht. Sein Flächeninhalt gibt die Integralkrümmung der Figur auf der gekrümmten Fläche an. Er ist positiv oder negativ zu nehmen, je nachdem, ob dieser Teil in Bezug auf seinen Rand ebenso liegt wie die ursprüngliche Figur in Bezug auf ihren Rand oder ob er umgekehrt liegt. Schneidet sich die Linie dagegen einmal oder mehrmals selbst, so ergibt sich eine komplizierte Figur, der man dennoch ebenso berechtigt einen bestimmten Flächeninhalt zuordnen kann wie Figuren ohne Knoten. So verstanden gibt dieser Flächeninhalt stets den richtigen Wert der Integralkrümmung an. Eine ausführlichere Darstellung %
